\chapter{Mensajes del sistema}
\label{ch:mensajes}
\cdtLabel{apendice:Mensajes}{}

%En esta sección se describen los mensajes utilizados en el prototipo actual del sistema.
Los mensajes se refieren a todos aquellos avisos que el sistema muestra al actor para comunicar la ocurrencia de algún evento tal como un error o una operación exitosa. Estos mensajes se pueden mostrar a través de diversos canales, por ejemplo, pantalla o correo electrónico.

Cuando un mensaje es recurrente se parametrizan sus elementos, por ejemplo los mensajes: ``Aún no existen registros de \textit{Unidades Académicas} en el sistema.'', ``Aún no existen registros de \textit{Programas Académicos} en el sistema.'', 
``Aún no existen registros de \textit{alumnos} en el sistema.'', tienen una estructura similar 
por lo que, con el objetivo de que el mensaje sea genérico y pueda utilizarse en todos los casos de uso que se considere necesario, se utilizan parámetros para definir el mensaje.\\

Los parámetros también se utilizan cuando la redacción del mensaje tiene datos que son ingresados por el actor o que dependen del resultado de la operación, por ejemplo: 
``La \textit{Unidad Académica ESCOM} ha sido \textit{modificada} exitosamente.''. En este caso la redacción se presenta parametrizada de la forma: ``$<DETERMINADO> <ENTIDAD> <VALOR>$ ha sido $<OPERACION>$ exitosamente.'' y los parámetros se describen de la siguiente forma:

\begin{itemize}
	\item DETERMINADO ENTIDAD: Es un artículo determinado más el nombre de la entidad sobre la cual se realizó la acción.
	\item VALOR: Es el valor asignado al atributo de la entidad, generalmente es el nombre o la clave.
	\item OPERACIÓN: Es la acción que el actor solicitó realizar.
\end{itemize}

En el ejemplo anterior se hace referencia a $<VALOR>$, es decir: \textit{ESCOM} es el \textbf{valor} de la \textbf{entidad} \textit{escuela}. Cada mensaje enlista los parámetros  que utiliza, sin embargo aquí se definen los más comunes a fin de simplificar la descripción de los mensajes:

\begin{description}
	\item [$<ARTICULO>$:] Se refiere a un {\em artículo} el cual puede ser DETERMINADO (El $\mid$ La $\mid$ Lo $\mid$ Los $\mid$ Las) o INDETERMINADO (Un $\mid$ Una $\mid$ 
	Uno $\mid$ Unos $\mid$Unas) se aplica generalmente sobre una ENTIDAD, ATRIBUTO o VALOR.
	
	\item [$<CAMPO>$:] Se refiere a un campo del formulario. Por lo regular es el nombre de un atributo en una entidad.
	
	\item [$<CONDICION>$:] Define una expresión booleana cuyo resultado deriva en {\em falso} o {\em verdadero} y suele ser la causa del mensaje.
	
	\item [$<DATO>$:] Es un sustantivo y generalmente se refiere a un atributo de una entidad descrito en el modelo estructural del negocio, por ejemplo: nombre de la unidad de aprendizaje, nombre del alumno, RFC del profesor, etc. %ATRIBUTO
	
	\item [$<ENTIDAD>$:] Es un sustantivo y generalmente se refiere a una entidad del modelo estructural del negocio, por ejemplo: programa académico, alumno, profesor, etc.
	\item [$<OPERACION>$:] Se refiere a una acción que se debe realizar sobre los datos de una o varias entidades. Por ejemplo: registrar, eliminar, actualizar, por mencionar algunos. Comúnmente la OPERACIÓN va concatenada con el sustantivo, por ejemplo: Registro de un nuevo beneficio, registro de una actividad, eliminar una tarea y demás.
	
	\item [$<VALOR>$:] Es un sustantivo concreto y generalmente se refiere a un valor en específico. Por ejemplo: \textit{Histología I} es un \textbf{valor} concreto de la \textbf{entidad} \textit{Unidad de Aprendizaje}.
	
\end{description}
%___________________________Plantilla_______________________________
%===========================  MSGX ==================================
%\begin{mensaje}{MSGX}{}{}
%	\item[Canal:] 
%	\item[Propósito:] 
%	\item[Redacción:]
%	\item[Parámetros:] 
%	\begin{itemize}
%		\item 
%	\end{itemize}
%	\item[Ejemplo:]  
%	\item[Referenciado por: ]
%\end{mensaje}

%===========================  MSG1 ==================================
\begin{mensaje}{MSG1}{Operación Exitosa}{Informativo }
	\item[Canal:] Sistema
	\item[Propósito:] Notificar al actor que la operación solicitada al sistema se llevó a cabo exitosamente.
	\item[Redacción:] La $<OPERACION>$ se llevo a cabo correctamente.
	\item[Parámetros:] OPERACION: Actividad que el actor debe realizar.
	\item[Ejemplo:] La reincripción se llevo a cabo correctamente.
\end{mensaje}

%============================== MSG2 =================================
\begin{mensaje}{MSG2}{Operación Fallida}{Error}
	\item[Canal:] Sistema
	\item[Propósito:] Notificar al actor que la operación solicitada al sistema no pudo llevarse a cabo.
	\item[Redacción:] Error, $<ERROR_ENCONTRADO>$.
	\item[Parámetros:] ERROR\_ENCONTRADO: Error que el sistema identificó al intentar realizar la operación.
	\item[Ejemplo:] Error, No hay conexión al sistema.(Base de Datos)
\end{mensaje}


%===========================  MSG3 ==================================
\begin{mensaje}{MSG3}{Elementos No Disponibles}{Error}
	\item[Canal:] Sistema
	\item[Propósito:] Notificar al actor que no existen elementos en el sistema por mostrar.
	\item[Redacción:] ¡No hay $<ELEMENTOS>$ para mostrarte!
	\item[Parámetros:] ELEMENTOS: Información de la cual se requiere para concluir un proceso
	\item[Ejemplo:] ¡No hay registros de aspirantes para mostrarte!
	
\end{mensaje}

%============================== MSG4 =================================
\begin{mensaje}{MSG4}{Elemento No Disponible}{Error}
	\item[Canal:] Sistema
	\item[Propósito:] Notificar al actor que el elemento que se desea ver no existe o no está disponible en el sistema.
	\item[Redacción:] ¡Este/Esta $<ELEMENTO>$ no esta disponible aún!
	\item[Parámetros:] ELEMENTO:Información de un registro en específico de la cual se requiere para concluir un proceso.
	\item[Ejemplo:] ¡Esta boleta no esta disponible aún!
	
\end{mensaje}

%============================== MSG5 =================================
\begin{mensaje}{MSG5}{Duplicado de Información}{Error}
	\item[Canal:] Sistema
	\item[Propósito:] Notificar al actor que ya existe un elemento con la misma información que desea ingresar.
	\item[Redacción:] ¡Ya existe un/a $<ELEMENTO>$ con la misma información!
	\item[Parámetros:] ELEMENTO: Parte de la información que esta duplicada
	\item[Ejemplo:]¡Ya existe una boleta con la misma información!
	
\end{mensaje}

%============================== MSG6 =================================
\begin{mensaje}{MSG6}{Falta dato obligatorio}{Error}
	\item[Canal:] Sistema
	\item[Propósito:] Notificar al actor la omisión de algún dato obligatorio por ingresar.
	\item[Redacción ] !Error¡ El campo $<DATO>$ es un campo obligatorio, favor de ingresar el campo faltante
	\item[Parámetros:] DATO: Información que se requiere para completar la operación.
	\item[Ejemplo:] !ERROR¡ El campo Referencia Bibliografíca es un campo obligatorio, favor de ingresar el campo faltante.
\end{mensaje}

%============================== MSG7 =================================
\begin{mensaje}{MSG7}{Formato de campo Incorrecto}{Error }
	\item[Canal:] Sistema
	\item[Propósito:] Notificar al actor que el dato ingresado en alguno de los campos del formulario no cumple con el tipo de dato definido en el diccionario de datos.
	\item[Redacción:] !Error¡ El campo $<DATO>$ ingresado es incorrecto, favor de ingresar un dato válido.
	\item[Parámetros:] DATO: Información que se requiere para completar la operación.
	\item[Ejemplo:] !ERROR¡ El campo Unidad de Aprendizaje ingresado es incorrecto, favor de ingresar un dato válido
\end{mensaje}


%===============================MSG8=============================

\begin{mensaje}{MSG8}{Operación por Atender}{Informativo }
	\item[Canal:] Correo / SMS / WhatsApp
	\item[Propósito:] Notificar al actor que existen operaciones que requieren de su atención o supervisión para llevarse a cabo.
	\item[Redacción:]Calmécac te notifica que hay $<ELEMENTOS>$ que requieren de tu atención.
	\item[Parámetros:]ELEMENTOS: Aquella información que el actor debe supervisar o atender.
	\item[Ejemplo:] Calmécac te notifica que hay registros de aspirantes que requieren de tu atención.
		
\end{mensaje}

%===============================MSG9=============================

\begin{mensaje}{MSG9}{Límite de Tiempo}{Informativo }
	\item[Canal:] Correo / SMS / WhatsApp
	\item[Propósito:] Notificar al actor que el límite de tiempo para que realice una operación está por agotarse.
	\item[Redacción:]¡ALERTA! Recuerda que debes realizar la $<OPERACION>$ antes del $<TIEMPO>$
	\item[Parámetros:]OPERACION: Tarea a realizar por el actor
	TIEMPO: Fecha para la cual el límite 
	\item[Ejemplo:] ¡ALERTA! Recuerda que debes realizar la reinscripción antes del 10 de Mayo del 2017.
		
\end{mensaje}


%===============================MSG10=============================
\begin{mensaje}{MSG10}{Cierre de Operación}{Informativo}
	\item[Canal:] Correo / SMS / WhatsApp
	\item[Propósito:] Notificar al actor que el límite de tiempo para que realice la operación se ha agotado y sugiere realizar algo para poder llevarla a cabo.
	\item[Redacción:] ¡ALERTA! El tiempo para realizar la $<OPERACION>$ se ha acabado. Calmécac te sugiere $<SUGERENCIA>$.
	\item[Parámetros:] 
	\begin{itemize}
		\item OPERACION: Tarea a realizar por el actor
		\item SUGERENCIA: Proceso u operación que el actor debe realizar para poder concluir una actividad.
	\end{itemize}
	\item[Ejemplo:] ¡ALERTA! El tiempo para realizar la reinscripción se ha acabado. Calmécac te sugiere ir al Departamento de Gestión Escolar para poder reinscribirte.
	
\end{mensaje}

%===========================  MSG11 ==================================
\begin{mensaje}{MSG11}{Recomendación para revisar configuración de servicios web}{Informativo}
	\item[Canal:] Correo / SMS / WhatsApp
	\item[Propósito:] Notificar al actor que la es necesario revisar la configuración actual de los servicios web, ya que se acerca un proceso de carga de datos al sistema.
	\item[Redacción:] Recomendación del Calmécac. La $<OPERACION>$ ha concluido, y se sugiere que se revise la configuración de los servicios web del Calmécac.
	\item[Parámetros:] OPERACION: Tarea realizada por el Politécnico.
	\item[Ejemplo:] Recomendación del Calmecac. La generación de citas ha concluido y se sugiere que se revise la configuración de los servicios web del Calmécac.
		
\end{mensaje}


%===========================  MSG12 ==================================
\begin{mensaje}{MSG12}{Confirmar operación sin cambios posteriores}{Confirmación}
	\item[Canal:] Sistema
	\item[Propósito:] Solicitar al actor la confirmación de una operación de impacto que no podrá ser revertida.
	\item[Redacción:] ¿Desea confirmar la finalización $<OPERACION>$? La operación no podrá ser revertida una vez confirmada.
	Calmécac te notifica que es necesaria tú aprobación para concluir la $<OPERACION>$
	\item[Parámetros:] OPERACION: Aquella operación que el actor debe confirmar.
	\item[Ejemplo:] ¿Desea confirmar la finalización del registro de aspirantes? La operación no podrá ser revertida una vez confirmada.
	Calmécac te notifica que es necesario tú aprobación para concluir el registro de aspirantes
	
\end{mensaje}

%===========================  MSG13 ==================================
\begin{mensaje}{MSG13}{Operación en Curso}{Informativo}
	\item[Canal:] Sistema
	\item[Propósito:] Notificar al actor que el sistema esta actualmente ejecutando una transacción.
	\item[Redacción:] Calmécac te notifica que se esta ejecutando la/el $<TRANSACCION>$
	\item[Parámetros:] TRASACCION: Aquella transaccion que se esta ejecutando actualmente
	\item[Ejemplo:] Calmécac te notifica que se esta ejecutando la generación de citas de reinscripción.
	    
\end{mensaje}


%===========================  MSG14 ==================================
\begin{mensaje}{MSG14}{No tienes notificaciones por atender}{Informativo}
	\item[Canal:] Correo / SMS / WhatsApp
	\item[Propósito:] Notificar al actor que no tiene mensajes o notificaciones que requieran de su atención.
	\item[Redacción:] Calmécac te notifica que no tienes mensajes o notificaciones pendientes que requieran de tu atención
\end{mensaje}

%===========================  MSG15 ==================================
\begin{mensaje}{MSG15}{No	 existe información en el sistema}{Error}
	\item[Canal:] Sistema
	\item[Propósito:] Notificar al actor que no se puede llevar a cabo la operación solicitada pues no existe la información suficiente en el sistema.
	\item[Redacción:] !Error! Falta información necesaria sobre $<INFORMACION>$ en el sistema para poder realizar la/el $<OPERACION>$ \\
	\item[Parámetros:] 
	\begin{Titemize}
		\Titem INFORMACIÓN: El tipo de información que hace falta en el sistema.
		\Titem OPERACION: Tarea a realizar por el actor
	\end{Titemize}
	\item[Ejemplo:] ¡Error! Falta información necesaria sobre Modalidad Educativa en el sistema para poder realizar el registro de Unidades de Aprendizaje
	\refIdElem{HR-UA-CU4}, \refIdElem{HR-UA-GG-CU1}
\end{mensaje}

%===========================  MSG16 ==================================
\begin{mensaje}{MSG16}{Condiciones necesarias para realizar operación}{Error}
	\item[Canal:] Sistema
	\item[Propósito:] Notificar al actor que no se puede realizar la operación solicitada debido a que no se cumplen las condiciones necesarias para hacerlo
	\item[Redacción:] No es posible $<OPERACION> <ARTICULO\_ENTIDAD>$ pues es necesario que $<CONDICIONES>$
	\item[Parámetros:] 
	\begin{itemize}
		\item OPERACION: Es la operación que el actor intentó llevar a cabo.
		\item ARTICULO\_ENTIDAD: Es la entidad sobre la cual se intentó realizar la operación acompañado del artículo de género correspondiente.
		\item CONDICIONES: Son las condiciones necesarias que no se satisfacieron.
	\end{itemize}
	\item[Ejemplo:] No es posible activar el plan de estudio pues es necesario que se especifiquen sus unidades de aprendizaje.
	
	
\end{mensaje}

%===========================  MSG17 ==================================
\begin{mensaje}{MSG17}{Eliminar elemento}{Confirmación}
	\item[Canal:] Sistema
	\item[Propósito:] Solicitar la confirmación del actor para la eliminación de un elemento.
	\item[Redacción:] ¿Desea eliminar $<ARTICULO>$ + $<ELEMENTO>$ + $<SELECCION>$?
	\item[Parámetros:] 
	\begin{itemize}
		\item ARTICULO: Es la parte de la oración que se ocupa de expresar el género (masculino/femenino).
		\item ELEMENTO: Es el elemento que se requiere eliminar.
		\item SELECCION:Dependiendo del género en la oración se ocupa, seleccionado(masculino)  ó seleccionada(femenino).
	\end{itemize}
	\item[Ejemplo:] ¿Desea eliminar el ciclo escolar seleccionado?
	
\end{mensaje}

%===========================  MSG18 ==================================
\begin{mensaje}{MSG18}{Relación de fechas incorrecta}{Error}
	\item[Canal:] Sistema
	\item[Propósito:] Notificar al actor que la relación entre dos fechas ingresadas no es la esperada.
	\item[Redacción:] La $<FECHA1>$ de $<ENTIDAD1>$ debe ser $<CONDICION>$ $<FECHA2>$ de $<ENTIDAD2>$
	\item[Parámetros:] 
	\begin{itemize}
		\item FECHA1: Fecha de una entidad.
		\item ENTIDAD1: Entidad que requiere de una fecha.
		\item CONDICION: La condición de relación que deben satisfacer FECHA1 Y FECHA2.
		\item FECHA2: Fecha de una entidad.
		\item ENTIDAD2: Entidad que requiere de una fecha.
	\end{itemize}
	\item[Ejemplo:] La fecha de fin de actividad debe ser menor que la fecha de fin de período escolar.
\end{mensaje}

%===========================  MSG19 ==================================
\begin{mensaje}{MSG19}{Relación de cantidades incorrecta}{Error}
	\item[Canal:] Sistema
	\item[Propósito:] Notificar al actor que la relación entre dos cantidades ingresadas no es la esperada.
	\item[Redacción:] La $<CANTIDAD1>$ debe ser $<CONDICION>$ $<CANTIDAD2>$.
	\item[Parámetros:] 
	\begin{itemize}
		\item CANTIDAD1: Cantidad de un elemento.
		\item CONDICION: La condición de relación que deben satisfacer CANTIDAD1 Y CANTIDAD2.
		\item CANTIDAD2: Cantidad de un elemento.
	\end{itemize}
	\item[Ejemplo:] La carga mínima debe ser menor que la carga media.
	
\end{mensaje}

%===========================  MSG20 ==================================
\begin{mensaje}{MSG20}{Relación de horas incorrecta}{Error}
	\item[Canal:] Sistema
	\item[Propósito:] Notificar al actor que la relación entre dos horas ingresadas no es la esperada.
	\item[Redacción:] Las(s) $<HORAS1>$ de $<ENTIDAD1>$ debe(n) ser $<CONDICION>$ $<HORAS2>$ de $<ENTIDAD2>$
	\item[Parámetros:] 
	\begin{itemize}
		\item HORAS1: Horas asociadas a una entidad.
		\item ENTIDAD1: Entidad que requiere de una asginación de horas.
		\item CONDICION: La condición de relación que deben satisfacer HORAS1 Y HORAS2.
		\item HORAS2: Horas asociadas a una entidad.
		\item ENTIDAD2: Entidad que requiere de una asginación de horas.
	\end{itemize}
	\item[Ejemplo:] Las horas de un contenido deben ser menores que las horas de la unidad de aprendizaje.
	
\end{mensaje}


%===========================  MSG21 ==================================
\begin{mensaje}{MSG21}{Elementos mínimos necesarios}{Error}
	\item[Canal:] Sistema
	\item[Propósito:] Notificar al actor que debe registrar al menos cierta cantidad de elementos para realizar la operación.
	\item[Redacción:] Es necesario $<OPERACION>$ al menos $<CANTIDAD>$ $<ENTIDAD>$.
	\item[Parámetros:] 
	\begin{itemize}
		\item OPERACIÓN: Operación que debe realizar el usuario. Puede ser registrar o seleccionar.
		\item CANTIDAD: La cantidad necesaria de elementos que el actor debe registrar.
		\item ENTIDAD: Es la entidad de la cuál se debe realizar la acción.
	\end{itemize}
	\item[Ejemplo:]  Es necesario registrar al menos un plan de estudio.
	
\end{mensaje}


%===========================  MSG22 ==================================
\begin{mensaje}{MSG22}{Notificación de operación.}{Informativo}
	\item[Canal:] Sistema
	\item[Propósito:] Notificar al actor la realización de una operación.
	\item[Redacción:] $<ARTICULO\_ENTIDAD>$ $<NOMBRE>$ ha sido $<OPERACION>$
	\item[Parámetros:] 
	\begin{itemize}
		\item ARTICULO\_ENTIDAD: Es la entidad a la que pertenece el elemento sobre el que se realizó la operación..
		\item NOMBRE: Es el nombre del elemento.
		\item OPERACIÓN: Es la operación que ha sido realizada.
	\end{itemize}
	\item[Ejemplo:] El período escolar 2017-2018/3-A1 ha sido eliminado.
	
\end{mensaje}

%===========================  MSG23 ==================================
\begin{mensaje}{MSG23}{No es posible realizar operación sobre elemento}{Informativo}
	\item[Canal:] Sistema
	\item[Propósito:] Notificar al actor que no es posible realizar una operación sobre un elemento debido a la existencia de otros elementos asociados a él.
	\item[Redacción:] No es posible $<OPERACION>$ $<ARTICULO\_ENTIDAD>$ $<NOMBRE>$ debido a que tiene $<ELEMENTOS\_ASOCIADOS>$ asociados.
	\item[Parámetros:] 
	\begin{itemize}
		\item OPERACIÓN: La operación que no se puede llevar a cabo.
		\item ARTICULO\_ENTIDAD: Es la entidad de la cuál se quiere eliminar un elemento.
		\item NOMBRE: Es el nombre del elemento.
		\item ELEMENTOS\_ASOCIADOS: Son los elementos que se encuentran asociados al elemento que se quiere eliminar.
	\end{itemize}
	\item[Ejemplo:] No es posible eliminar el edificio Edificio Inteligente debido a que tiene espacios asociados. 
	
\end{mensaje}

%===========================  MSG24 ==================================
\begin{mensaje}{MSG24}{Confirmación de eliminación en cadena}{Confirmación}
	\item[Canal:] Sistema
	\item[Propósito:] Solicitar al actor la confirmación para eliminar un elemento que tiene como consecuencia la elimación de otros elementos.
	\item[Redacción:] Al eliminar $<ARTICULO\_ENTIDAD>$ $<NOMBRE>$ se eliminarán $<ARTICULO\_ENTIDAD\_ASOCIADA>$ que se listan a continuación, ¿desea continuar?
	\item[Parámetros:] 
	\begin{itemize}
		\item ARTICULO\_ENTIDAD: Es la entidad de la cuál se eliminará el elemento.
		\item NOMBRE: Es el identificador del elemento.
		\item ARTICULO\_ENTIDAD\_ASOCIADA: Es la entidad de las cuál se eliminarán elementos debido a la eliminación de NOMBRE.
	\end{itemize}
	\item[Ejemplo:]  Al eliminar el periodo escolar se eliminarán las estructuras académicas que se listan a continuación, ¿desea continuar?
	
\end{mensaje}

%===========================  MSG25 ==================================
\begin{mensaje}{MSG25}{Confirmación de operación}{Confirmación}
	\item[Canal:] Sistema
	\item[Propósito:] Solicitar al actor la confirmación para llevar a cabo una operación.
	\item[Redacción:] ¿Desea $<OPERACION>$ $<ARTICULO\_ENTIDAD>$ $<NOMBRE>$?
	\item[Parámetros:] 
	\begin{itemize}
		\item OPERACIÓN: Es la operación que el actor desea realizar.
		\item ARTICULO\_ENTIDAD: Es la entidad de la cuál se eliminará el elemento.
		\item NOMBRE: Es el identificador del elemento.
	\end{itemize}
	\item[Ejemplo:] ¿Desea editar el plan de estudio Plan 2009?
\end{mensaje}

%===========================  MSG58 ==================================
\begin{mensaje}{MSG58}{Operación fuera de periodo}{Error}
	\item[Canal:] Sistema
	\item[Propósito:] Notificar al actor que la operación que desea realizar está fuera del periodo permitido
	\item[Redacción:] No es posible realizar $<ARTICULO\_OPERACION>$ debido al que la fecha actual está fuera del periodo permitido para realizarla.
	\item[Parámetros:] 
	\begin{itemize}
		\item ARTICULO\_OPERACION: La operación que el actor desea llevar a cabo.
	\end{itemize}
	\item[Ejemplo:] No es posible realizar el envío de la solicitud debido a que la fecha actual está fuera del periodo permitido para realizarla.
	\refIdElem{HR-PR-CU1}
\end{mensaje}


